\documentclass[a4paper]{article}

\usepackage{graphicx}
\usepackage{amsmath,amssymb}
\usepackage{minted}
\usepackage{multirow}
\usepackage{float}
\usepackage{algorithm}
\usepackage{algpseudocode}
\usepackage{todonotes}
\usepackage{forest}
\usepackage{xstring}
\usepackage{xcolor}
\usepackage[most]{tcolorbox}
\usepackage{xparse}
\usepackage{natbib}
\usepackage{adjustbox}
\usepackage[active,tightpage]{preview}

\PreviewBorder=4pt

\title{Visual mathematics rendering and sizing fix}
\date{2026-09-16}
\author{Mohammad Rahmani}

\begin{document}
\maketitle

\section{Fix}
Visual mode now preloads both inline and display mathematics in the same
LuaLaTeX batch before the QTextDocument projection is constructed. Inline
constructs such as \verb|$X$|, \verb|$x_i$| and \verb|\(z\)| therefore use
rendered math images instead of remaining as blue source-backed chips.

Display formulas are normalized against the active Visual editor font metrics
before insertion. Small one-line formulas are enlarged to a readable size,
while multi-row constructs such as \texttt{align} and nested \texttt{cases}
reserve proportional vertical space. Width remains capped by the Visual
viewport so large equations are never forced beyond the page.

The normal path still performs one LuaLaTeX compilation for all uncached math
in a source load. Late inline tokens deliberately do not start individual TeX
processes, preserving the CPU behavior of the batch renderer.

\end{document}
